% GENERATED FILE. Edit canonical lessons, then run npm run generate:books.
% canonical-source-hash: fnv1a64:9036eb7e89cd3918
% canonical-lessons: ZH-C402-chengke, ZH-C402-houche, ZH-C402-gaotie, ZH-C402-huancheng, ZH-C402-chexiang

\chapter{At the Station}
\label{ch:zh-at-the-station}
\hlchaptermodality{402}

\begin{chapteropening}
\textbf{By the end of this chapter:} \emph{I can talk about passengers, the waiting room, fast trains, changing trains and the carriages.}

Complete the last lesson of chapter 402: I can talk about passengers, the waiting room, fast trains, changing trains and the carriages.
\end{chapteropening}

\section[chéngkè]{chéngkè — passenger}

\label{lesson:ZH-C402-chengke}



\begin{quote}
Before the new one: say the Mandarin for to get old, then the Mandarin for everything.
\end{quote}

\subsection*{You'll want to know: chéngkè}
\emph{chéngkè} — \textquotedblleft{}passenger\textquotedblright{}.

\begin{quote}
\emph{huǒchēzhàn lǐ yǒu yìxiē chéngkè} — There are a few passengers in the station.
\end{quote}

Its characters come later; for now it is a word for the ear.

\begin{cousinweb}
\emph{chéngkè} is built from \emph{chéng} 'to ride' and \emph{kè} 'a guest', both inherited from Middle Chinese, \emph{zying} and \emph{khaek} (Guangyun, in Baxter's transcription), which glosses the first 'to drive', 'to mount' and 'to prevail'.
\end{cousinweb}

\begin{grammarlens}[title={What you've built}]
One more word for getting around.
\end{grammarlens}

\subsection*{Your turn}
\begin{itemize}
\raggedright
  \item \emph{Say it:} \emph{chéngkè}
  \item \emph{Say it:} \emph{chéngkè}, once more
  \item \emph{Say it:} \emph{huǒchēzhàn lǐ yǒu yìxiē chéngkè} — There are a few passengers in the station.
  \item \emph{Recall:} say the Mandarin for to get old, then the Mandarin for everything, then say \emph{chéngkè} again
\end{itemize}

\begin{tcolorbox}[breakable,colback=teal!4,colframe=teal!35!black,title={Before you move on}]
What does chéngkè mean? (\textquotedblleft{}Passenger\textquotedblright{}.) Say it once more.
\end{tcolorbox}
\section[\texorpdfstring{\zh{候车} (hòuchē)}{候车 (hòuchē)}]{\zh{候车} (hòuchē) — to wait for a train or bus}

\label{lesson:ZH-C402-houche}



\begin{quote}
Before the new one: say the Mandarin for everything, then the Mandarin for a passenger.
\end{quote}

\subsection*{You'll want to know: \zh{候车}}
\textbf{\zh{候车}} — \emph{hòuchē} — \textquotedblleft{}to wait for a train or bus\textquotedblright{}.

\begin{quote}
\textbf{\zh{我们在车站候车}} \emph{wǒmen zài chēzhàn hòuchē} — We're waiting for the train at the station.
\end{quote}

\begin{cousinweb}
\textbf{\zh{候车}} (\emph{hòuchē}) is built from \zh{候} 'to wait for' and \zh{车} 'vehicle', both inherited from Middle Chinese, \emph{huwH} and \emph{tsyhae} (Guangyun, in Baxter's transcription), which glosses the first 'to wait upon, to watch for'.
\end{cousinweb}

\begin{grammarlens}[title={What you've built}]
One more word for getting around.
\end{grammarlens}

\subsection*{Your turn}
\begin{itemize}
\raggedright
  \item \emph{Say it:} \emph{hòuchē}
  \item \emph{Say it:} \emph{hòuchē}, once more
  \item \emph{Say it:} \emph{wǒmen zài chēzhàn hòuchē} — We're waiting for the train at the station.
  \item \emph{Recall:} say the Mandarin for everything, then the Mandarin for a passenger, then say \emph{hòuchē} again
\end{itemize}

\begin{tcolorbox}[breakable,colback=teal!4,colframe=teal!35!black,title={Before you move on}]
What does \zh{候车} mean? (\textquotedblleft{}To wait for a train or bus\textquotedblright{}.) Say it once more.
\end{tcolorbox}
\section[gāotiě]{gāotiě — high-speed railway}

\label{lesson:ZH-C402-gaotie}



\begin{quote}
Before the new one: say the Mandarin for a passenger, then the Mandarin for to wait for a train.
\end{quote}

\subsection*{You'll want to know: gāotiě}
\emph{gāotiě} — \textquotedblleft{}high-speed railway\textquotedblright{}.

\begin{quote}
\emph{wǒmen zuò gāotiě qù Běijīng} — We're taking the high-speed train to Beijing.
\end{quote}

Its characters come later; for now it is a word for the ear.

\begin{cousinweb}
\emph{gāotiě} is built from \zh{高} 'high' and \emph{tiě} 'iron', both inherited from Middle Chinese, \emph{kaw} and \emph{thet} (Guangyun, in Baxter's transcription), which glosses the second, after an older dictionary, 'black metal'. It is a short form of a longer name meaning a high-speed iron road.
\end{cousinweb}

\begin{grammarlens}[title={What you've built}]
One more word for getting around.
\end{grammarlens}

\subsection*{Your turn}
\begin{itemize}
\raggedright
  \item \emph{Say it:} \emph{gāotiě}
  \item \emph{Say it:} \emph{gāotiě}, once more
  \item \emph{Say it:} \emph{wǒmen zuò gāotiě qù Běijīng} — We're taking the high-speed train to Beijing.
  \item \emph{Recall:} say the Mandarin for a passenger, then the Mandarin for to wait for a train, then say \emph{gāotiě} again
\end{itemize}

\begin{tcolorbox}[breakable,colback=teal!4,colframe=teal!35!black,title={Before you move on}]
What does gāotiě mean? (\textquotedblleft{}High-speed railway\textquotedblright{}.) Say it once more.
\end{tcolorbox}
\section[huànchéng]{huànchéng — to change trains or lines}

\label{lesson:ZH-C402-huancheng}



\begin{quote}
Before the new one: say the Mandarin for to wait for a train, then the Mandarin for the high-speed railway.
\end{quote}

\subsection*{You'll want to know: huànchéng}
\emph{huànchéng} — \textquotedblleft{}to change trains or lines\textquotedblright{}.

\begin{quote}
\emph{wǒmen zài zhèr huànchéng} — We change trains here.
\end{quote}

Its characters come later; for now it is a word for the ear.

\begin{cousinweb}
\emph{huànchéng} is built from \emph{huàn} 'to exchange' and \emph{chéng} 'to ride', both inherited from Middle Chinese, \emph{hwanH} and \emph{zying} (Guangyun, in Baxter's transcription), which glosses the first 'to exchange'.
\end{cousinweb}

\begin{grammarlens}[title={What you've built}]
One more word for getting around.
\end{grammarlens}

\subsection*{Your turn}
\begin{itemize}
\raggedright
  \item \emph{Say it:} \emph{huànchéng}
  \item \emph{Say it:} \emph{huànchéng}, once more
  \item \emph{Say it:} \emph{wǒmen zài zhèr huànchéng} — We change trains here.
  \item \emph{Recall:} say the Mandarin for to wait for a train, then the Mandarin for the high-speed railway, then say \emph{huànchéng} again
\end{itemize}

\begin{tcolorbox}[breakable,colback=teal!4,colframe=teal!35!black,title={Before you move on}]
What does huànchéng mean? (\textquotedblleft{}To change trains or lines\textquotedblright{}.) Say it once more.
\end{tcolorbox}
\section[chēxiāng]{chēxiāng — railway carriage}

\label{lesson:ZH-C402-chexiang}



\begin{quote}
Before the new one: say the Mandarin for the high-speed railway, then the Mandarin for to change trains.
\end{quote}

\subsection*{You'll want to know: chēxiāng}
\emph{chēxiāng} — \textquotedblleft{}railway carriage\textquotedblright{}.

\begin{quote}
\emph{wǒmen zuò zài sān hào chēxiāng} — We're sitting in carriage three.
\end{quote}

Its characters come later; for now it is a word for the ear.

\begin{cousinweb}
\emph{chēxiāng} is built from \zh{车} 'vehicle' and \emph{xiāng} 'a side room', both inherited from Middle Chinese, \emph{tsyhae} and \emph{sjang} (Guangyun, in Baxter's transcription), which glosses the second 'a gallery' and the rooms east and west of a hall.
\end{cousinweb}

\begin{grammarlens}[title={What you've built}]
One more word for getting around.
\end{grammarlens}

\subsection*{Your turn}
\begin{itemize}
\raggedright
  \item \emph{Say it:} \emph{chēxiāng}
  \item \emph{Say it:} \emph{chēxiāng}, once more
  \item \emph{Say it:} \emph{wǒmen zuò zài sān hào chēxiāng} — We're sitting in carriage three.
  \item \emph{Recall:} say the Mandarin for the high-speed railway, then the Mandarin for to change trains, then say \emph{chēxiāng} again
\end{itemize}

\begin{tcolorbox}[breakable,colback=teal!4,colframe=teal!35!black,title={Before you move on}]
What does chēxiāng mean? (\textquotedblleft{}Railway carriage\textquotedblright{}.) Say it once more.
\end{tcolorbox}
